
\needspace{4\baselineskip}
\color{uniblue}{\subsection{Контроль синхронизма (КС)}}
\color{black}

\begin{enumerate}[label=\arabic{section}.\arabic{subsection}.\arabic*, labelsep=4pt, leftmargin=0pt, itemindent=57pt]

\item
Контроль синхронизма применяется, как правило, на выключателях линий с двухсторонним питанием, шиносоединительных и секционных выключателях. Данная функция предназначена для соединения двух независимых систем на параллельную работу через выключатель, при этом напряжения по двум системам синхронны, а угол между напряжениями не превышает заданного значения. Функция может применяться совместно с АПВ или при оперативном включении выключателя.

\item
В устройстве контроль синхронизма представлен функцией «КС».

\item
Основные характеристики функции:
\begin{itemize}
    \item погрешность функции по разности частот не более 20 мГц;
    \item погрешность функции по разности фаз не более $5^\circ$;
    \item погрешность срабатывания функции по скорости изменения частоты не более 150 мГц; 
    \item время срабатывания и возврата функции по разности частот не более 60 мс;
    \item время срабатывания ИО скорости изменения частоты не более 200 мс.
\end{itemize}

\item
Проверка синхронизма выполняется путем сравнения одного из линейных напряжений основного ТН с соответствующим напряжением ТН присоединения.

\item
Положением уставки SGF2 выбирается синхронизируемое напряжение, величины которого используются для выполнения расчета ($U_{AB}$, $U_{BC}$, $U_{CA}$). Если к устройству подается только одно линейное напряжение от ТН присоединения, то положением уставки SGF2 следует выбрать именно то напряжение, которое подведено к устройству. Если же подводятся напряжения всех трех фаз, то допускается применять в качестве синхронизируемого напряжения любое линейное напряжение.

\item
Если при одинаковых первичных напряжениях сдвиги фаз вторичных напряжений, подаваемых на устройство, сдвинуты относительно друг друга на определенный угол, то возможно скомпенсировать данный сдвиг при помощи уставки «Ф\_сдвига».

\item
Разрешение на оперативное включение выключателя «КС: Разреш.опер.вкл.В» формируется при выполнении следующих условий:
\begin{itemize}
    \item оперативный вывод КС (постоянный разрешающий сигнал);
    \item по приходу запроса на включение с контролем синхронизма от КП или УВ и выполнения условий по синхронизму (с контролем симметрии напряжений от КННш, КННп);
    \item по приходу запроса на включение с контролем синхронизма от КП или УВ и выполнения условий по напряжениям (нет Uш и нет Uп, нет Uш и есть Uп, нет Uп и есть Uш).
\end{itemize}

\item
Разрешение на автоматическое включение выключателя «КС: Разреш.авт.вкл. В» формируется по приходу запроса на включение с контролем синхронизма от АПВ и выполнения условий по синхронизму (с контролем симметрии напряжений от КННш, КННп).

\item Функция формирует следующие основные сигналы:
\begin{itemize}
    \item «КС: ОС» (при выполнении условий ожидания синхронизма);
    \item «КС: УС» (при выполнении условий улавливания синхронизма) при заданной уставки «Предусмотрено» программного переключателя SGF3 «Улавл\_синхр».
\end{itemize}

\item
Ожидание синхронизма (ОС) применяется при частоте скольжения после аварийного отключения в режиме максимальной нагрузки не более 0,4 Гц.

\item Логика ожидания синхронизма (ОС) формирует следующие вспомогательные сигналы:
\begin{itemize}
    \item «КС: Нет условий по U» при недопустимой разности модулей напряжений на основном ТН и ТН присоединения;
    \item «КС: Нет условий по Фос» при недопустимой разности углов между векторами напряжений основного ТН и ТН присоединения для ОС;
    \item «КС: Нет условий по Fос» при недопустимой разности частот напряжений основного ТН и ТН присоединения для ОС.
\end{itemize}   

\item
Улавливание синхронизма (УС) применяется при частоте скольжения после аварийного отключения в режиме максимальной нагрузки до 2 Гц. В устройстве реализован алгоритм, позволяющий включить выключатель при наименьшем отклонении углов между векторами напряжений на основном ТН и ТН присоединения. Данный алгоритм учитывает текущую скорость и ускорение скольжения УС с постоянным временем опережения.

\item Логика улавливания синхронизма (УС) дополнительно формирует следующие вспомогательные сигналы (при введенной логике улавливания синхронизма):
\begin{itemize}
    \item «КС: Нет условий по фус» при недопустимой разности углов между векторами напряжений основного ТН и ТН присоединения для УС;
    \item «КС: Блокировка по dFdt»  при недопустимой разности скорости изменения частот напряжений основного ТН и ТН присоединения;
    \item «КС: Нет условий по Fус» при недопустимой разности частот напряжений основного ТН и ТН присоединения для УС;
\end{itemize}  

\item
Пуск ИО по разности напряжений происходит при снижении модуля разности напряжения ниже уставки. Разность напряжений рассчитывается по формуле:
% ФОРМУЛА
\begin{equation*}
   dU= \lvert U_{\text{ш}} -  U_{\text{п}}   \rvert,
\end{equation*}
\begin{eqexpl}[25mm]
    \item{$U_{\text{ш}}$} напряжение основного ТН;
    \item{$U_{\text{п}}$} напряжение ТН присоединения.
\end{eqexpl}
Сигнал «КС: Нет условий по U» формируется при наличии сигнала общего пуска и отсутствии срабатывания ИО по разности напряжения.

\item
Пуск ИО по разности фаз для ожидания синхронизма происходит при снижении модуля разности фаз ниже уставки. Разность фаз рассчитывается по формуле:
% ФОРМУЛА
\begin{equation*}
   d\Phi= \lvert \Phi_{\text{ш}} - \Phi_{\text{п}} + \Phi_{\text{сдвига}} \rvert ,
\end{equation*}
\begin{eqexpl}[25mm]
    \item{$\Phi_{\text{ш}}$} фаза напряжение основного ТН;
    \item{$\Phi_{\text{п}}$} фаза напряжение ТН присоединения;
    \item {$\Phi_{\text{сдвига}}$} угол сдвига (компенсирует сдвиг напряжения основного ТН относительно ТН присоединения при синхронной работе).
\end{eqexpl}
Сигнал «КС: Нет условий по Фос» формируется при наличии сигнала общего пуска и отсутствии срабатывания ИО по разности фаз для ожидания синхронизма.

\item
Пуск ИО по разности фаз для улавливания синхронизма происходит при снижении величины фазы включения (с учетом времени опережения) ниже уставки. Величина фазы включения рассчитывается по формуле:
% ФОРМУЛА
\begin{equation*}
   \Phi_{\text{вкл.}}= d\Phi +  360 \cdot dF \cdot T_{\text{опер.}},
\end{equation*}
\begin{eqexpl}[25mm]
    \item{$d\Phi$} разность фаз векторов напряжений;
    \item{$dF$} разность частот векторов напряжения;
    \item{$T_{\text{опер.}}$} уставка функции.
\end{eqexpl}
Чтобы включение выключателя было выполнено с минимальной разницей по фазе, команда на включение будет подаваться с учетом времени работы выключателя, задаваемой уставкой «Tопер».
Сигнал «КС: Нет условий по Фус» формируется при введенной в работу уставке SGF3, наличии сигнала общего пуска и отсутствии срабатывания ИО по разности фаз для улавливания синхронизма.

\item
Пуск ИО по разности частот для ожидания синхронизма происходит при снижении модуля разности частот сети основного ТН и ТН присоединения ниже уставки. Возврат ИО будет выполняться при превышении разности частот выше значения (dFср + dFвоз). Разность частот рассчитывается по формуле:

% ФОРМУЛА
\begin{equation*}
    dF= \lvert F_{\text{ш}} - F_{\text{п}} \rvert ,
 \end{equation*}
 \begin{eqexpl}[25mm]
    \item{$F_{\text{ш}}$} частота сети основного ТН;
    \item{$F_{\text{п}}$} частота сети ТН присоединения.
 \end{eqexpl}
 Сигнал «КС: Нет условий по Fос» формируется при наличии сигнала общего пуска и отсутствии срабатывания ИО по разности частот для ожидания синхронизма.

\item
Пуск ИО по разности частот для улавливания синхронизма происходит при снижении модуля  разности частот сети основного ТН и ТН присоединения энергообъектов ниже уставки. Возврат ИО будет выполняться при превышении разности частот выше значения (dFмакс\_ср + dFмакс\_воз). Сигнал «КС: Нет условий по Fус» формируется при введенной в работу уставке SGF3, наличии сигнала общего пуска и отсутствии срабатывания ИО по разности частот для улавливания синхронизма.
\item
Пуск ИО по скорости изменения частот происходит при выполнении одного из условий:
\begin{itemize}
\item превышении модуля скорости изменения частоты на основном ТН выше уставки «dFdt\_ср»;
\item превышении модуля скорости изменения частоты на ТН присоединения выше уставки «dFdt\_ср»;
\item превышении модуля разности скоростей изменения частоты на основном ТН и ТН присоединения выше уставки «dFdt\_ср».
\end{itemize}

\item
Сигнал «КС: Блокировка по dF/dt» формируется при введенной в работу уставке SGF3, наличии сигнала общего пуска и срабатывания ИО по скорости изменения частот.

\item
Ввод функции в работу на этапе параметрирования устройства осуществляется программным переключателем SGF1 «Ввод\_функции». Информация о введенной функции отображается выходным сигналом «КС: Ввод».

\item
Функция может быть выведена из работы оперативно путем активации сигнала «ОВ КС», что характеризуется активным состоянием сигнала «КС: Оперативный вывод» на выходе функции.

% ПРЕДУПРЕЖДЕНИЕ !
%\needspace{4\baselineskip}
\par
\noindent
\begin{tabularx}{\linewidth}{p{2em} X}
  \begin{minipage}[c][3\baselineskip][b]{2em}
    \centering\textcolor{unired}{\Huge\faTriangleExclamation}
  \end{minipage} &
  \hyphenpenalty=0 \textbf{Оперативный вывод КС дает безусловное разрешение на оперативное включение выключателя без контроля синхронизма, что может привести к ошибочному несинхронному включению оборудования.}
\end{tabularx}
\vskip 2mm
%%%%%%%%%%%%%%%%%%%

\item
Алгоритм работы функции выполнен в соответствии с рисунком \ref{ks:img1}. В таблице \ref{ks:tbl1} приведены параметры, необходимые для настройки функции.

% УСЛОВНОЕ ВЫРАЖЕНИЕ
\ifboolexpr{bool {isVV}}
{
\vspace{1mm}
\begin{figure}[H]
\centering
\includegraphics[width=0.9\textwidth,height=0.9\textheight,keepaspectratio]{\fbpath/126. КС/_latex/img/ksvv.pdf}
\captionof{figure}{Функциональная схема логики КС}\label{ks:img1}
\end{figure}
}
{
\vspace{1mm}
\begin{figure}[H]
\centering
\includegraphics[width=0.9\textwidth,height=0.9\textheight,keepaspectratio]{\fbpath/126. КС/_latex/img/ks.pdf}
\captionof{figure}{Функциональная схема логики КС}\label{ks:img1}
\end{figure}
}

\end{enumerate}
